% BÀI TOÁN HÌNH HỌC - BAI_HINH_26 (TRÍCH ĐỀ 26)
\begin{ex}
Từ điểm M nằm ngoài đường tròn (O) kẻ hai tiếp tuyến MA, MB tới (O) (A, B là các tiếp điểm). Gọi H là giao điểm của AB với MO; E và F là các giao điểm của đường thẳng MO với đường tròn (O) (với ME < MF).
	\begin{enumerate}
		\item[1.] Chứng minh bốn điểm M, A, O, B cùng thuộc một đường tròn.
		\item[2.] Chứng minh MO $\perp$ AB và $HE \cdot HF = HM \cdot HO$.
		\item[3.] Kẻ đường kính BP của đường tròn (O). Đường thẳng MP cắt đường tròn (O) tại điểm N (N khác P). Gọi I là trung điểm của đoạn thẳng MH. Chứng minh ba điểm A, N, I thẳng hàng.
	\end{enumerate}
\end{ex}
